% RESEARCH_DRAFT_SHA256: 15df82bd66eb770972a5355a838e17fa007bc010e00ba45b03ed3310db3abaad
% Layout and source-reference wrappers expanded for the web renderer.
\documentclass[11pt,a4paper]{article}
\usepackage[margin=25mm]{geometry}
\usepackage{fontspec}
\setmainfont{lmroman10-regular.otf}[BoldFont=lmroman10-bold.otf,
 ItalicFont=lmroman10-italic.otf,BoldItalicFont=lmroman10-bolditalic.otf]
\usepackage{amsmath,amssymb,mathtools}
\usepackage{unicode-math}
\setmathfont{latinmodern-math.otf}
\AtBeginDocument{\let\setminus\smallsetminus}
\usepackage{xurl,hyperref}
\hypersetup{colorlinks=true,urlcolor=blue}
\setcounter{secnumdepth}{0}
\setlength{\parindent}{0pt}
\setlength{\parskip}{5pt}
\setlength{\emergencystretch}{3em}
\begin{document}
\section*{Research-attempt provenance}
This standalone research attempt retains the complete original definition body below,
including its historical source and preparation notes. Those notes describe the
earlier catalogue preparation; the new attempt and its verification are identified
separately after that body. No original ID, catalogue status, or source statement is changed.

\section*{Hamiltonian Tetrahedralizations}\label{5500029}
{\small UnsolvedMath ID 5500029 · AMR-054-0029 · Computer Science · AMR Open Problem Lists}
\subsection{Definitions and mathematical statement}
Let $P\subset\mathbb R^3$ be a convex three-dimensional polytope. A
tetrahedralization $\mathcal T$ partitions $P$ into finitely many
tetrahedra meeting face-to-face, with vertices drawn only from the
vertices of $P$. Subdivision of a nonsimplicial boundary face by
diagonals is allowed; insertion of new vertices is not.

The dual graph $G_{\mathcal T}$ has one vertex for each tetrahedron,
with two vertices adjacent precisely when the corresponding
tetrahedra share a triangular face. It has a Hamiltonian path if its
tetrahedra can be ordered $T_1,\ldots,T_m$ so that every tetrahedron
occurs once and $T_i,T_{i+1}$ share a face for each $i<m$.

Does every convex $P$ admit a tetrahedralization with this property?
The subdivision itself may be chosen to make the path possible;
the claim is not that every tetrahedralization works. The path is
in the dual graph, not a path through the original polytope vertices.

Allowing new Steiner points gives a different problem and is known
to permit additional constructions. The source's unrestricted
vertex-preserving existence question is the one retained here.
No bound on the number of tetrahedra beyond finiteness is part of
the main assertion.
\subsection{Short English statement}
Can every convex three-dimensional polytope be tetrahedralized, using
its existing vertices, so that all tetrahedra can be visited once in
an order crossing a shared face at each step?
\subsection{Sources}
\begin{itemize}
\item[S1] ULAM.AI and the UnsolvedMath Contributors, supplied problems.json, record 5500029 (AMR-054-0029), AMR Open Problem Lists. Catalogue identity and source transcription; CC BY 4.0.
\url{https://huggingface.co/datasets/ulamai/UnsolvedMath}.
\item[S2] Open Problems Project - Computational Geometry, Problem 29. Original source indexed by AMR-054-0029.
\url{https://topp.openproblem.net/p29}.
\end{itemize}

\section{Literature and context check}
The consulted TOPP Problem 29 page labels the question open and records
obstructions for particular pulling triangulations. The supplied prior narrative
asserting a general positive solution is internally inconsistent with that source;
this attempt does not adopt it. The argument below is independent of that
literature and tests a specific shelling-based shortcut.
\section{Research attempt: Shellability does not supply a Hamiltonian dual path}
A necessary condition for a finite graph $G$ to possess a Hamiltonian path is
\[
 \#\operatorname{components}(G-S)\le |S|+1
 \quad\hbox{for every vertex set }S.
\]
Indeed, deleting $|S|$ vertices from a path leaves at most $|S|+1$ nonempty
path pieces. Each component of $G-S$ contains at least one such piece because
the path visited every vertex; graph edges may merge pieces but cannot increase
the number of components.

Now start with a tetrahedron $T$ and stack one tetrahedron on each of three
distinct boundary faces, with no further stacking. This is geometrically
realizable as a convex polytope with a vertex-preserving tetrahedralization.
Here is a justification of the geometric step often omitted from the graph
picture. At any chosen triangular facet of a convex polytope, its barycenter
satisfies strictly all other facet inequalities. Moving it a sufficiently small
positive distance along the outward normal keeps those finitely many strict
inequalities and violates only the chosen facet inequality. The convex hull
with this new vertex adds exactly the pyramid over that facet; the visible
boundary consists of that facet alone. The other two original facets remain
facets. Apply this operation successively on the three selected original facets.
Each new vertex is extreme, so all tetrahedra use only vertices of the resulting
convex polytope.

The resulting four-tetrahedron dual graph is $K_{1,3}$: each new tetrahedron
shares a triangle with $T$, while two new ones meet, if at all, along an edge
or vertex of $T$. The order consisting of $T$ followed by its three attachments
is a shelling in the direct sense that each attachment meets the preceding
union in a single triangular facet. Yet the dual has no Hamiltonian path:
deleting its center leaves three components, violating the bound $3\le2$.

Thus shelling, connectivity, and even convex realizability of every intermediate
stacking step do not turn an arbitrary traversal into a face-adjacent
Hamiltonian traversal. A depth-first walk would return through the central
tetrahedron repeatedly. The problem asks for existence of some triangulation,
however; the demonstrated failure of this triangulation leaves alternative
tetrahedralizations entirely available. Any repair must modify the subdivision
to eliminate such dual separation obstructions, rather than merely reorder the
given shelling.
\section{Verification and adversarial audit}
The construction uses three separate facet stackings, and the strict
inequality argument verifies convexity and extreme vertices without placing
Steiner vertices in the interior. The dual's four vertices and three edges were
counted explicitly. One or two attached tetrahedra instead have path duals,
confirming the smallest branching obstruction in this family.
\paragraph{Audit corrections.} The prior unsupported solved narrative was rejected. Failure of one shellable triangulation was not confused with failure of the existential original target.
\paragraph{External computation.} None; symbolic arguments and boundary checks only.
\section{Scope and remaining gaps}
There is no proof that every convex polytope has a triangulation whose dual
avoids all separating-set obstructions, and avoiding those obstructions alone
would still be only necessary for Hamiltonicity. No negative example for all
tetrahedralizations of one polytope is claimed.
\paragraph{Integration review.} Codex (GPT-6; exact serving revision not exposed) reviewed the written arguments and expanded layout and source-reference wrappers. Repository metadata and the production web renderer were checked using Python 3.13 and JavaScript with MathJax 3.2.2. This model-assisted review is not independent human proof verification; no PDF compilation or formal proof checking is asserted.

\section{Completion Estimate}
\noindent\textbf{COMPLETION: 0\%}
The calculations below the preserved source statement investigate a concrete route,
but establish no substantiated new progress toward the entire original target.
Known special cases and failed sufficient conditions are not credited as a solution.
\section{References}
The original source references and catalogue attribution are preserved in the
Sources section above. The following URLs were actually consulted for this attempt;
retention of an additional original reference does not claim fresh consultation of it.
\begin{itemize}
\item \url{https://topp.openproblem.net/p29} (consulted 2026-10-11). 
\end{itemize}
\section{Model and contributor disclosure}
OpenAI gpt-6-astra, requested reasoning configuration ultra, generated this research
attempt on 2026-10-11. The exact serving revision is not exposed. Codex GPT-6
performed orchestration and repository-format review; its exact serving revision
is also not exposed. The mathematical route, boundary checks and final adversarial
self-audit followed contribution prompts A then B in
\texttt{docs/contribution\_prompts.md}, after reading \texttt{CONTRIBUTING.md}.
Python 3.14.2 assembled the standalone documents; any mathematical computation is
specified in Verification. No theorem prover was used. No human mathematical
authorship, substantial human proof correction, or independent human proof checking
is asserted. Model-assisted review, source retrieval, and structural checks do not
establish correctness or novelty. Full private reasoning transcripts are unavailable
as reproducibility artifacts; the complete arguments, inputs and outcomes actually
used are stated here or in the companion batch audit.
\section{Final status}
Unresolved. This is a research attempt with an explicit failed route or restricted
calculation, not a full solution, not a current-status certification, and not a claim
that the elementary material is new. The catalogue's original source-status label
is preserved independently of this attempt status.
\end{document}
